\section{Advantages and limitations (G3)}\label{sec:discussion}
Table~\ref{tab:advlim} summarises the advantages and limitations that follow from the decision logic and from the audit; the empirical columns of Section~\ref{sec:results} qualify them. Three cross-cutting limitations recur. (1) The V1/V2 gates select on \emph{training} $\R$, so nothing in the gate protects against a model that overfits and then extrapolates badly; V6 is the only variant selecting on held-out edge data. (2) Ensemble weighting is defined three different ways across V1--V3, so ``ensemble'' is not one component. (3) Several variants execute LLM-generated code with \texttt{exec}, with only a restricted namespace as isolation.
\begin{table}[t]\centering\scriptsize\rtab{t11_adv_lim}\caption{Advantages and limitations per variant.}\label{tab:advlim}\end{table}
\paragraph{Reliability discrepancy.} One DeFi case (Portfolio Variance) reports $\R=1.000$ in-sample and $\rfar=-882.9$ at seed 42 for the same variant. Figure~\ref{fig:trainfar} shows train versus far $\R$ over all runs; the experiment \texttt{hybrid\_exp1b} is its multi-seed sweep. We do not average this case away.
\begin{figure}[t]\centering\rfig[0.55]{f6_train_vs_far}\caption{Train vs.\ far $\R$ (symlog axis).}\label{fig:trainfar}\end{figure}
